\lesson{3}{Sep 14 2021 Tue (16:54:18)}{Solving Radical Equations}{Unit 1}

\subsubsection{Solving Radical Equations}
\label{sub_sub_sec:solving_radical_equations}

Let's go over how to solve this radical equation: $\sqrt{(5x + 4)} = 7$:

\begin{enumerate}
  \label{enum:solving_radical_equations}
  \item Isolate the Radical:
    When solving equations, the goal is always to isolate the variable to find
    its value. The goal is the same with radical equations. So, we must make the
    radical isolated first. But in this case, we already have the radical
    isolated.
  \item Square both Sides: Do the opposite operation of taking the square root
    and square both sides of the equation.
      \begin{equation}
        \label{eqt:solving_radical_equations_square_both_sides_1}
        \begin{split}
          \sqrt{(5x + 4)} &= 7 \\
          \sqrt{(5x + 4)}^{2} &= 7^{2} \\
          5x + 4 &= 49 \\
        \end{split}
      .\end{equation}
    Once the radical is gone, we just simply solve for the variable.
  \item Solve for the Variable: To solve a two-step equation, remember that you
    must use the reverse order of operations. \\
    (SADMEP: \textbf{S}ubtraction, \textbf{A}ddition, \textbf{D}ivision,
    \textbf{M}ultiplication, \textbf{E}xponents, \textbf{P}arentheses)
      \begin{equation}
        \label{eqt:solving_radical_equations_solve_for_variable_2}
        \begin{split}
          5x + 4  &= 49 \\
          5x      &= 45 \\
          \frac{5x}{5} &= \frac{45}{5} \\
          \boxed{x = 9}
        \end{split}
      .\end{equation}
  \item Check your Work: When solving radical equations, you must always check
    your work. Even if you do all of the work perfectly, you may not find the
    correct solution. So, let's plug the value of $9$ into the original equation
    to see if it works.
      \begin{equation}
        \label{eqt:solving_radical_equations_check_your_work_3}
        \begin{split}
          \sqrt{5x + 4} &= 7 \\
          \sqrt{5(9) + 4} &= 7 \\
          \sqrt{45 + 4} &= 7 \\
          \sqrt{49} &= 7 \\
          \boxed{7 = 7}
        \end{split}
      .\end{equation}
\end{enumerate}

Let's take a look at a radical equation that doesn't have any solutions:

\begin{example}
  \label{exm:solving_radical_equations_no_solutions}

  \begin{equation}
    \label{eqt:solving_radical_equations_no_solutions_1}
    \begin{split}
      \sqrt{x - 3} + 4 &= 1 \\
      \sqrt{x - 3} &= -3 \\
      \sqrt{x - 3}^{2} &= -3^{2} \\
      x - 3 &= 9 \\
      \boxed{x = 12}
    \end{split}
  .\end{equation}

  Now, let's check our work:

  \begin{equation}
    \label{eqt:solving_radical_equations_no_solutions_2}
    \begin{split}
      \sqrt{x - 3} + 4 &= 1 \\
      \sqrt{12 - 3} + 4 &= 1 \\
      \sqrt{9} + 4 &= 1 \\
      3 + 4 &= 1 \\
      7 \ne 1 \\
    \end{split}
  .\end{equation}

  This statement is false. This is known as a \textbf{extraneous solution}.
  (\ref{def:extraneous_solution})
\end{example}

\begin{definition}[Extraneous Solution]
  \label{def:extraneous_solution}

  An \textbf{Extraneous Solution} is a solution to an equation that doesn't fit
  the requirements of the original equation.
\end{definition}

% subsubsection solving_radical_equations (end)

\subsubsection{Solving Radical Equations with Variables Outside the Radical}
\label{sub_sub_sec:solving_radical_equations_with_variables_outside_the_radical}

Before we start to solve the equation below
(\ref{eqt:solving_radical_equations_with_variables_outside_the_radical}), take
a look at the following property (\ref{prop:zero_product_property})

\begin{property}[Zero Product Property]
  \label{prop:zero_product_property}

  If $a \times b = 0$, then $a = 0$ and/or $b = 0$

  Let's say we are given the following polynomial: $(x + 8)(x - 4) = 0$. Here's
  how we would solve it. First, we take each item within the parentheses and set
  that equal to $0$. For example:

  \[ (x + 8) = 0 \quad (x - 4) = 0 \].

  Now, let's find the zeros of the polynomial.

  \[ x = -8 \quad x = 4 \].
\end{property}

Let's solve: $\sqrt{x + 9} - 7 = x$:

\begin{equation}
  \label{eqt:solving_radical_equations_with_variables_outside_the_radical}
  \begin{split}
    \sqrt{x + 9} - 7 &= x \\
    \sqrt{x + 9} &= x + 7 \quad \textrm{Isolate the variable} \\
    (\sqrt{x + 9})^{2} &= (x + 7)^{2} \quad \textrm{Square both sides} \\
    x + 9 &= (x + 7)(x + 7) \\
    x + 9 &= x^{2} + 14x + 49 \quad \textrm{Solve for the variable} \\
    9 &= x^{2} + 13x + 49 \\
    0 &= x^{2} + 13x + 40 \quad \textrm{Factor} \\
    0 &= (x + 5)(x + 8) \quad \textrm{Apply the zero product property
      (\ref{prop:zero_product_property}}) \\
    x + 5 &= 0 \quad \textrm{AND} \quad x + 8 = 0 \\
    \boxed{x = -5, -8} \quad \textrm{Check your Work} \\
    \sqrt{x + 9} - 7 &= x \quad \sqrt{x + 9} - 7 = x \\
    \sqrt{-5 + 9} - 7 &= -5 \quad \sqrt{-8 + 9} - 7 = -8 \\
    \sqrt{4} - 7 &= -5 \quad \sqrt{1} - 7 = -8 \\
    2 - 7 &= -5 \quad 1 - 7 = -8 \\
    -5 &= -5 \quad -6 \ne -8 \\
  \end{split}
.\end{equation}

$x = -5$ is the only solution and $x = 8$ is extraneous.

% subsubsection solving_radical_equations_with_variables_outside_the_radical (end)

\newpage
